\section*{Project contract}
Implement \texttt{run\_vm} for the instruction set defined in this chapter. Pass
the full program-byte address in \texttt{a0}, its length in \texttt{a1}, a state
descriptor address in \texttt{a2}, and a virtual-instruction budget in \texttt{a3}.
The program has zero through 64 bytes inside a supplied 96-byte arena and may start
at an unaligned byte address. The budget is zero through 64. The descriptor is
eight-byte aligned and names a word-aligned operand-stack subview with capacity zero
through eight words. Program, descriptor and stack allocations are disjoint.

Each call starts at program byte zero with an empty logical operand stack and zero
retired guest instructions. Descriptor output fields from a previous call are ignored.
The call validates the complete encoded image, validates every encoded branch target,
and then executes until \texttt{HALT}, a runtime fault, or budget exhaustion. A
budget counts successfully retired guest instructions; it excludes validation work
and does not count native MIPS instructions. Exhaustion is checked before the next
fetch. A zero budget still performs both validation phases.

Return the status in \texttt{v0} and the retired guest count in \texttt{v1}. Publish
PC, depth, retired count, status and fault offset in the selected descriptor. The PC
is a byte offset. Status zero means budget exhaustion and status one means halt.
Negative statuses identify the failures specified below. A successful halt advances
PC past its opcode and counts as one retired instruction. A failing instruction does
not retire; retain its PC and the live operand values from before that instruction.
Validation failures report PC, depth and retired count zero, with the offending
instruction's offset in the separate fault field.

The program is the only bytecode input. Descriptor fields and the selected stack
capacity may be read. Caller bytes outside descriptor output fields and stack capacity
are preserved, including temporary writes. Program bytes, other arenas, capacity,
stack pointer and guards are read-only. Inactive stack cells within capacity are
unspecified workspace; live values occupy the first depth words. Submission-owned
constants and private MIPS stack storage are permitted. Preserve full \texttt{s0} through \texttt{s7}, \texttt{fp} and \texttt{sp};
return through \texttt{ra} with no console output. Equivalent decoding and dispatch
implementations are accepted.
